\editionPage[\pdfbookmark[0]{Fiches outils}{bm-outils-debut-parent}]{outils-debut}{Fiche outil}{Calculs et algèbre}{À consulter au fil des chapitres.}
\begin{blocEdition}{outils-debut}

\editionEtape{Fractions et relatifs}Pour additionner des fractions, utiliser un même dénominateur. Pour multiplier, multiplier les numérateurs entre eux et les dénominateurs entre eux. Pour diviser par une fraction non nulle, multiplier par son inverse. Avant un produit de relatifs, contrôler le signe.
\editionEtape{Puissances et racines}Une puissance positive est un produit répété. $10^{-n}=1/10^n$ ; l’écriture scientifique positive est $a\times10^n$ avec $1\le a<10$. Pour $b\ge0$, $\sqrt b$ est non négatif et son carré vaut $b$.
\editionEtape{Distributivité}$k(a+b)=ka+kb$. Développer va du produit vers la somme ; factoriser fait le chemin inverse. Réduire uniquement des termes de même nature.
\editionEtape{Équation}Développer et réduire, appliquer les mêmes opérations aux deux membres, isoler l’inconnue. Vérifier ensuite dans l’équation initiale et dans le contexte.
\editionRetenir{Une égalité exacte utilise $=$. Une valeur arrondie utilise $\approx$. Conserver le résultat exact ou la valeur complète de la calculatrice jusqu’au dernier calcul.}
\end{blocEdition}
\editionPage{outils-2}{Fiche outil}{Conversions et formules}{À consulter au fil des chapitres.}
\begin{blocEdition}{outils-2}

\editionEtape{Temps et vitesses}$1$ h $=60$ min $=3600$ s. $v=d/t$, $d=vt$, $t=d/v$. De km/h à m/s : diviser par 3,6. Dans l’autre sens : multiplier par 3,6.
\editionEtape{Volumes et débits}$1$ m$^3=1000$ L ; $1$ L $=1$ dm$^3$ ; $1$ mL $=1$ cm$^3$. $q=V/t$. Pour convertir un débit, traiter séparément le volume et l’unité de temps.
\editionEtape{Aires}Rectangle : $L\ell$. Triangle : base $\times$ hauteur associée divisée par 2. Disque : $\pi r^2$. Circonférence : $2\pi r$. Une hauteur de triangle est perpendiculaire à la droite de sa base.
\editionEtape{Volumes}Pavé : $L\ell h$. Prisme et cylindre : aire de base $\times h$. Pyramide et cône : aire de base $\times h/3$. Pour le cône : $\pi r^2h/3$.
\editionEtape{Agrandir ou réduire}Un rapport positif $k$ agit sur les longueurs par $k$, sur les aires par $k^2$ et sur les volumes par $k^3$. Les angles sont conservés.
\editionRetenir{Avant de multiplier deux mesures, choisir des unités compatibles. Une longueur en cm et une aire en cm$^2$ donnent un volume en cm$^3$.}
\end{blocEdition}
\editionPage{outils-3}{Fiche outil}{Choisir un outil géométrique}{À consulter au fil des chapitres.}
\begin{blocEdition}{outils-3}

\editionEtape{Une longueur dans un triangle rectangle}Deux côtés connus : Pythagore. Un angle aigu et un côté connu : cosinus si l’on dispose de l’adjacent ou de l’hypoténuse. Écrire le rapport avant de choisir une opération.
\editionEtape{Démontrer un angle droit}Comparer le carré du plus long côté à la somme des deux autres carrés. Si les deux calculs sont égaux, utiliser la réciproque de Pythagore et nommer le sommet opposé au plus long côté.
\editionEtape{Des parallèles et des triangles emboîtés}Thalès pour calculer à partir d’un parallélisme connu. Sa réciproque pour prouver un parallélisme à partir de rapports égaux. Vérifier alignements, ordre des points et correspondance.
\editionEtape{Comparer deux triangles}Chercher CCC, CAC ou ACA. Respecter les sommets correspondants, puis transférer une longueur ou un angle. Une translation conserve les mesures et fournit des parallélogrammes.
\editionRetenir{Une preuve cite ses données et sa propriété. Un codage issu de l’énoncé est une donnée ; une mesure approximative du dessin est seulement un indice pour chercher.}
\end{blocEdition}
\editionPage{outils-4}{Fiche outil}{Données, hasard et programmation}{À consulter au fil des chapitres.}
\begin{blocEdition}{outils-4}

\editionEtape{Données}Fréquence = effectif de la catégorie / effectif total. Moyenne = somme des valeurs répétées / effectif total. Pour une médiane, ranger les valeurs et trouver les rangs centraux. Angle d’un secteur = fréquence $\times360^\circ$.
\editionEtape{Hasard}Définir l’expérience et les issues. Vérifier l’équiprobabilité avant de compter favorable / total. Pour le contraire : $1-P(A)$. Dans un modèle uniforme sur une cible, comparer les aires. Ne pas confondre modèle de probabilité et fréquence observée.
\editionEtape{Graphiques}Lire titre, axes, unités et graduations. Pour une lecture inverse, parcourir toutes les intersections de l’horizontale avec le tracé. Une droite doit passer par l’origine pour représenter une proportionnalité.
\editionEtape{Scripts}Repérer le déclencheur ; initialiser ; suivre l’ordre des blocs ; noter les variables ; compter les tours ; tester les cas limites. Pour un dessin, préciser départ, orientation et stylo. Dans un tableur, une référence relative change quand on recopie une formule.
\editionRetenir{Après correction, refaire une question sans aide et expliquer l’idée qui l’a résolue. Le but est de retrouver la méthode, pas de reconnaître le résultat.}
\end{blocEdition}
